% ED26-SRC-063-0103 continuation, original 2004 chapter §§2.3--7, pp.166--243.
% Reader fragment: follows src063-0103; it does not introduce another chapter.
\section{The categorical deduction apparatus}
\label{sec:p13-boger-deduction}

Boger's deduction system has four categorical sentence kinds, the contradictory
pairs \(a/o\) and \(e/i\), the universal contrary pair \(a/e\), three conversion
rules, fourteen two-premise syllogism rules, and direct and indirect deduction.
These are ingredients of a reconstruction of Aristotle's underlying logic,
not fourteen axiomatic conditional sentences. An interpreted argument, its
pattern, and the deduction making its validity evident remain different objects
\parencite[pp.~166--178]{Boger2004UnderlyingLogic}.

The syntax here remains \emph{predicate first}: \(AaB\) says that \(A\) belongs
to every \(B\). The letters are places for homogeneous terms, not individual
variables. Johnson's \(Aab\), with subject \(a\) first, must not be copied into
these derivations. Nor does the present apparatus contain Johnson's separate
assertion and rejection judgments. Boger's inconcludence procedure classifies
premise-pair patterns; it is not a rejection calculus obtained by writing a
reversed assertion sign.

\begin{hbdefinition}{def:p13-boger-deduction-domain}{Scope of the categorical derivations}
In the ordinary syllogistic cases developed below, terms have an existent
subject matter and \emph{some} is indeterminate: at least one, possibly all.
Universal affirmations carry the existence needed for conversion per accidens.
Universal privatives can be weakened to partial privatives only with the
relevant subject's existence. The contradictory of \(AaB\) is \(AoB\), and
that of \(AeB\) is \(AiB\), with the same term order. The reverse assignments
give the other two contradictories.
\end{hbdefinition}

This is a stated local policy, not erasure of the earlier nonexistence discussion.
In particular, one cannot combine vacuous universal affirmations about empty
subjects with conversion per accidens. The reduction of Darapti below also
uses the incompatibility of \(AaB\) and \(AeB\) for an existent \(B\).
Boger explicitly relies on nonempty subject matter in the contrasted-instance
and weakening examples \parencite[pp.~193--197]{Boger2004UnderlyingLogic}.

\subsection{Formation before inference}
\label{sec:p13-boger-deduction-formation}

The conversion-premise formation rules CPFR1--3 take exactly one categorical
sentence, with different terms, and close that formation account. CAFR1--4
give one premise and one conclusion, exchange the two term positions, and
retain the categorical constant for \emph{simple} conversion. The change from
\(a\) to \(i\) is the expressly exceptional conversion per accidens; CAFR2
must not be read as forbidding that exception.

SPFR1--8 take exactly two categorical sentences involving exactly three
different terms, never repeating a term within one sentence. One term occurs
in both premises. Each premise may have any of the four constants. The shared
term is the middle. Its three positions yield the three figures. SAFR1--8
then fix one categorical conclusion made from the extremes, rather than from
the middle. The conclusion's predicate is the major extreme; its subject is
the minor extreme. These are formation restrictions, not guarantees of validity
\parencite[pp.~167--175]{Boger2004UnderlyingLogic}.

\begin{tabularx}{\textwidth}{@{}l X X@{}}
\toprule
Figure & Premise-pair pattern & Conclusion place\\
\midrule
First, PMS & \(PxM,MyS\) & \(PzS\)\\
Second, MPS & \(MxP,MyS\) & \(PzS\)\\
Third, PSM & \(PxM,SyM\) & \(PzS\)\\
\bottomrule
\end{tabularx}

Here \(x,y,z\) are metalinguistic places for \(a,e,i,o\), not quantified
individuals. Premise order does not change consequence. Its ordered display
does fix the major/minor bookkeeping and makes an enumeration systematic.
There are \(3\cdot4^2=48\) ordered premise-pair patterns in these three figures,
and \(3\cdot4^3=192\) patterns when a conclusion kind is specified. Boger's
fourteen rules select concludent pairs and their characteristic conclusions;
they do not assert that every alternative conclusion from those pairs is valid.
Weakened conclusions remain separate consequences.

\subsection{Conversions and the fourteen rules}
\label{sec:p13-boger-deduction-rules}

CR1 changes \(AeB\) to \(BeA\); CR2 changes \(AaB\) to \(BiA\);
CR3 changes \(AiB\) to \(BiA\). There is no simple \(a\)-conversion and
no \(o\)-conversion. Boger's account of CR1 uses Aristotle's selection of an
instance from an assumed positive converse. His table~12's notation \(C=A\)
is an admitted simplification, not identity of classes or a new identity axiom
(note~70). Tables~13--14 then assume \(BeA\), use CR1 to get \(AeB\), and
conflict respectively with \(AaB\) and \(AiB\). The first conflict requires
existent \(B\); the second is a contradictory pair. Thus even the conversion
argument has an existence dependency, not merely a permutation of letters
\parencite[pp.~180--183]{Boger2004UnderlyingLogic}.

\begin{hblemma}{thm:p13-boger-conversions}{proved=annex:p13-boger-conversions}{The three conversions under the ordinary policy}
For homogeneous categorical terms, interpreted in the ordinary existent-subject
scope, \(AeB\) entails \(BeA\), \(AaB\) entails \(BiA\), and \(AiB\)
entails \(BiA\). Simple conversion of \(AaB\) to \(BaA\), or of \(AoB\)
to \(BoA\), is not licensed by these meanings.
\end{hblemma}

The annex expands the property-bearing cases and the source's indirect
arguments. This supplies a local justification of the admitted rules; it
does not attribute an unbounded quantifier calculus or a class-identity
axiom to Aristotle.

For the failure of \(o\)-conversion, Aristotle's terms are man and animal.
Man does not belong to every animal, while animal belongs to every man.
The proposed converse therefore has a true premise and false conclusion.
This \emph{method of fact} is different from the two-instantiation method
subsequently used for premise-pair inconcludence.

The literal rule displays are collected in
\cref{sheet:p13-boger-deduction}. The following list preserves the source's
term orientation and states the rule inventory without treating it as a proof.

\begin{tabularx}{\textwidth}{@{}l l X l@{}}
\toprule
Source rule & Name & Premises & Conclusion\\
\midrule
SR1 & Barbara & \(AaB,BaC\) & \(AaC\)\\
SR2 & Celarent & \(AeB,BaC\) & \(AeC\)\\
SR3 & Darii & \(AaB,BiC\) & \(AiC\)\\
SR4 & Ferio & \(AeB,BiC\) & \(AoC\)\\
SR5 & Cesare & \(MeN,MaX\) & \(NeX\)\\
SR6 & Camestres & \(MaN,MeX\) & \(NeX\)\\
SR7 & Festino & \(MeN,MiX\) & \(NoX\)\\
SR8 & Baroco & \(MaN,MoX\) & \(NoX\)\\
SR9 & Darapti & \(PaS,RaS\) & \(PiR\)\\
SR10 & Felapton & \(PeS,RaS\) & \(PoR\)\\
SR11 & Disamis & \(PiS,RaS\) & \(PiR\)\\
SR12 & Datisi & \(PaS,RiS\) & \(PiR\)\\
SR13 & Bocardo & \(PoS,RaS\) & \(PoR\)\\
SR14 & Ferison & \(PeS,RiS\) & \(PoR\)\\
\bottomrule
\end{tabularx}

This inventory agrees with the legible SR1--14 statements and table~20.
It is neither the traditional twenty-four-mood four-figure presentation nor
Johnson's axiomatic assertion/rejection system.

\section{Completion: how the additional rules are established}
\label{sec:p13-boger-completion}

For Boger, a first-figure \emph{teleios sullogismos} is a valid argument whose
validity is evident without an additional chain. An \emph{atel\={e}s} argument
is already valid, but requires a deduction to make that validity evident.
Completion does not confer truth on premises or manufacture validity by
recognition. At the metalogical level, completion makes an additional pattern's
panvalidity evident by using the conversions and the four first-figure patterns
as rules \parencite[pp.~183--188]{Boger2004UnderlyingLogic}.

\begin{hbdefinition}{def:p13-boger-deduction-bookkeeping}{Direct and indirect chains}
For this reader's bookkeeping, \(\Gamma\vdash_{\mathrm{full}}c\) means that
there is a finite categorical deduction of \(c\) from the premise set
\(\Gamma\). An available premise may be repeated; an admitted conversion or
syllogism rule may be applied to earlier lines. An indirect subdeduction adds
the categorical contradictory \(c^{*}\) temporarily, derives incompatible
sentences using the stated opposition policy, and discharges that assumption
to conclude \(c\). This star and this derivability notation are editorial
metalinguistic aids, not new categorical operators or ancient judgment signs.
\end{hbdefinition}

\calculussheet{sheet:p13-boger-deduction}{Boger's categorical deduction apparatus}{Predicate-first categorical patterns; finite premise-to-conclusion chains; categorical opposition and assumption discharge; no individual-variable equality calculus.}
\hbrule{rule:p13-boger-conv-e}{CR1 e-conversion}{AeB}{BeA}{Homogeneous different terms; ordinary syllogistic scope.}
\hbrule{rule:p13-boger-conv-a}{CR2 conversion per accidens}{AaB}{BiA}{The universal affirmation carries existent subject B.}
\hbrule{rule:p13-boger-conv-i}{CR3 i-conversion}{AiB}{BiA}{Indeterminate positive particular.}
\hbrule{rule:p13-boger-barbara}{SR1 Barbara}{AaB\quad BaC}{AaC}{Predicate-first term order.}
\hbrule{rule:p13-boger-celarent}{SR2 Celarent}{AeB\quad BaC}{AeC}{Predicate-first term order.}
\hbrule{rule:p13-boger-darii}{SR3 Darii}{AaB\quad BiC}{AiC}{Indeterminate some.}
\hbrule{rule:p13-boger-ferio}{SR4 Ferio}{AeB\quad BiC}{AoC}{Indeterminate some.}
\hbrule{rule:p13-boger-cesare}{SR5 Cesare}{MeN\quad MaX}{NeX}{Established using CR1 and SR2.}
\hbrule{rule:p13-boger-camestres}{SR6 Camestres}{MaN\quad MeX}{NeX}{Established using CR1 and SR2.}
\hbrule{rule:p13-boger-festino}{SR7 Festino}{MeN\quad MiX}{NoX}{Established using CR1 and SR4; reduced-system proof differs.}
\hbrule{rule:p13-boger-baroco}{SR8 Baroco}{MaN\quad MoX}{NoX}{Indirect completion using SR1.}
\hbrule{rule:p13-boger-darapti}{SR9 Darapti}{PaS\quad RaS}{PiR}{Existence enters through CR2.}
\hbrule{rule:p13-boger-felapton}{SR10 Felapton}{PeS\quad RaS}{PoR}{Existence enters through CR2.}
\hbrule{rule:p13-boger-disamis}{SR11 Disamis}{PiS\quad RaS}{PiR}{Established using CR3 and SR3.}
\hbrule{rule:p13-boger-datisi}{SR12 Datisi}{PaS\quad RiS}{PiR}{Established using CR3 and SR3.}
\hbrule{rule:p13-boger-bocardo}{SR13 Bocardo}{PoS\quad RaS}{PoR}{Indirect completion using SR1.}
\hbrule{rule:p13-boger-ferison}{SR14 Ferison}{PeS\quad RiS}{PoR}{Established using CR3 and SR4.}
\hbrule{rule:p13-boger-repetition}{Repetition}{\Gamma\vdash_{\mathrm{full}}c}{\Gamma\vdash_{\mathrm{full}}c}{Editorial context notation; repeat an available line.}
\hbrule{rule:p13-boger-indirect}{Categorical indirect deduction}{\Gamma,c^{*}\vdash q\quad\Gamma,c^{*}\vdash r}{\Gamma\vdash c}{q and r are incompatible under the stated opposition and existence policy; discharge only c-star.}


The direct chain for Camestres does more than mention Celarent. From
\(MaN,MeX\), convert the second premise to \(XeM\); Celarent gives
\(XeN\); conversion gives \(NeX\). The last conversion restores the prescribed
major/minor orientation. For Baroco, instead assume \(NaX\), the contradictory
of the intended \(NoX\). Barbara with \(MaN\) gives \(MaX\), contradicting
the actual premise \(MoX\). Discharging \(NaX\) gives \(NoX\).
These are respectively direct and indirect completion, not modern
first-order proofs substituted for the source's categorical steps.

\begin{hbtheorem}{thm:p13-boger-ten-completions}{proved=annex:p13-boger-ten-completions}{All ten additional rules are obtainable}
With CR1--3, the four first-figure rules, repetition, and categorical indirect
deduction under \cref{def:p13-boger-deduction-domain}, each of SR5--14 has a
finite deduction from its two displayed premises to its displayed conclusion.
\end{hbtheorem}

The annex supplies all ten chains, not just Camestres and Baroco.
They follow the original table~19: Cesare and Camestres use Celarent;
Festino, Felapton, and Ferison use Ferio; Darapti, Disamis, and Datisi
use Darii; Baroco and Bocardo use Barbara under a discharged contradictory
assumption. A full derivation of a pattern establishes its availability as a
rule; a list of rule names alone does not.

\section{Deciding inconcludence without conflating procedures}
\label{sec:p13-boger-inconcludence}

\subsection{Ordinary contrasted instances}
\label{sec:p13-boger-ordinary-contrast}

Boger follows Ross in naming Aristotle's procedure \emph{contrasted instances}.
For the first-figure pair \(PaM,MeS\), take animal, man, horse: both premises
are true, and animal belongs to every horse. Then take animal, man, stone:
both premises are true, and animal belongs to no stone. The displayed
conclusions are \emph{both true}, in different substitutions. They therefore
do not themselves constitute the modern true-premises/false-conclusion test
\parencite[pp.~188--192]{Boger2004UnderlyingLogic}.

The procedure's force concerns the shared premise-pair pattern. The second
instance rules out a necessary affirmative result; the first rules out a
necessary privative result. With existent subjects and indeterminate
particulars, those two exclusions cover all four conclusion kinds. Switching
the conclusions gives the modern counterarguments: the stone instance makes
the universal affirmative false, and the horse instance makes the universal
privative false. This adaptation explains a relationship between the methods;
it does not retroactively change which procedure the source describes.

\begin{hbtheorem}{thm:p13-boger-contrast-four}{proved=annex:p13-boger-contrast-four}{What a contrasted pair excludes}
Fix a categorical premise-pair pattern and the standard conclusion order
\(PzS\). Suppose one admissible substitution makes both premises and \(PaS\)
true, and another makes both premises and \(PeS\) true, with an existent
subject in each. Assume the ordinary opposition policy of
\cref{def:p13-boger-deduction-domain}. No conclusion kind \(a,e,i,o\) in
that fixed order follows necessarily from the premise-pair pattern.
\end{hbtheorem}

This is a scoped semantic statement about substitutions. It does not prove
that every interpreted argument with extra information about the chosen terms
is invalid. Boger's stronger language of paninvalidity belongs to the
elementary pattern analysis; additional topical axioms change the premise
information and must be recorded rather than smuggled into the pattern.

\subsection{Modified selection and the indeterminate route}
\label{sec:p13-boger-modified-contrast}

For \(PaM,MoS\) and \(PeM,MoS\), the partial privative premise obstructs
the ordinary construction. Aristotle's modified example selects swan and snow
from white things to which man does not belong. Animal is predicated of all
of the selected swans and none of the selected snow; with inanimate in the
privative major premise the two results reverse. Boger emphasizes selection
from an existent subject matter and the determinate use of \emph{some} here.
The selected part must not silently become a claim about \emph{all} white
things. His tables~23--24 are abbreviated representations of that selection
\parencite[pp.~192--194]{Boger2004UnderlyingLogic}.

A different route uses \emph{some} indeterminately. A universal denial entails
the corresponding partial denial when its subject exists. Thus an already
inconcludent stronger pair can witness inconcludence of its weakened pair.
For example, \(MeP,MeS\) is strengthened relative to \(MeP,MoS\).
Aristotle's second-figure discussion cannot obtain the affirmative contrast
with a determinate \(MoS\): if \(PaS\) and \(MeP\) hold, Celarent gives
\(MeS\), conflicting with the positive part of a determinate reading. He
therefore changes to the indeterminate route; he does not simply omit the
missing instance \parencite[pp.~195--197]{Boger2004UnderlyingLogic}.

\begin{hbtheorem}{thm:p13-boger-weakening-inconcludence}{proved=annex:p13-boger-weakening-inconcludence}{Nonconsequence survives premise weakening}
Let \(\Gamma_{\mathrm{s}}\) and \(\Gamma_{\mathrm{w}}\) be stronger and weaker
premise-pair patterns over the same terms and conclusion order. Suppose every
admissible substitution satisfying \(\Gamma_{\mathrm{s}}\) also satisfies
\(\Gamma_{\mathrm{w}}\). If each of the four candidate conclusions has a
counterinstance for \(\Gamma_{\mathrm{s}}\), it has one for
\(\Gamma_{\mathrm{w}}\). In particular, replacing an \(e\) premise by its
indeterminate \(o\) counterpart preserves this nonconsequence conclusion
when the subject exists.
\end{hbtheorem}

The implication between premise conditions is essential. A determinate
\emph{some but not all} privative requires a positive as well as a negative
case and is not entailed by a universal denial. Weakening cannot be justified
just by calling a sentence ``partial.''

\subsection{The exhaustive premise-pair classification and its boundary}
\label{sec:p13-boger-pair-catalogue}

The source catalogues the unselected pairs in tables~26--28 and the selected
pairs in SR1--14. In this table a two-letter code records the major and minor
categorical constants, in that order, within the figure already fixed above.
The two columns partition its sixteen premise-pair combinations in each figure.

\begin{tabularx}{\textwidth}{@{}l l X@{}}
\toprule
Figure & Concludent pairs & Inconcludent pairs in the source catalogue\\
\midrule
First & \(aa,ea,ai,ei\) & \(ae,ee,ia,oa,ie,oe,ao,eo,ii,oo,io,oi\)\\
Second & \(ea,ae,ei,ao\) & \(aa,ee,oa,ie,eo,ai,oe,ia,ii,oo,io,oi\)\\
Third & \(aa,ea,ia,ai,oa,ei\) & \(ae,ee,ao,ie,oe,eo,ii,oo,io,oi\)\\
\bottomrule
\end{tabularx}

This is the forty-eight-pair source classification, not forty-eight independently
proved reader theorems. A concludent pair still admits invalid candidate
conclusions. Barbara's pair, for example, establishes \(PaS\) (and a weakened
\(PiS\)), not \(PeS\) or \(PoS\). Darii does not establish a universal
affirmation simply because its particular affirmation follows. The alternative
conclusion therefore has to be examined even when the pair is concludent
\parencite[pp.~202--207]{Boger2004UnderlyingLogic}.

Boger calls elementary patterns non-neutrovalid: form, within the stipulated
categorical information, determines their validity. His condition~I excludes
the two opposed universal outcomes under true-premise substitutions; condition~II
allows both outcomes in different such substitutions. His epistemic principles
distinguish a deduction establishing panvalidity from contrasted instances
establishing inconcludence. These are claims in his reconstruction, not a
general theorem that any two unrelated true sentences are logically contrary
\parencite[pp.~191--202]{Boger2004UnderlyingLogic}.

Changing conclusion order is another operation. Simple \(e\) and \(i\)
conversion preserves an obtained conclusion. An \(a\) conclusion only gives
an \(i\) converse; an \(o\) conclusion has no conversion rule. This explains
why reversing Barbara's conclusion does not give a universal converse.
Boger discusses a putative fourth figure through conversions and premise
transpositions, and distinguishes the resulting traditional names from
Aristotle's three-figure formation scheme. His discussion of an \(a\)-converse
per accidens leading to Bramantip is itself a warning against treating every
change of orientation as preserving the original conclusion strength.

\section{Reduction: a smaller sufficient syllogism-rule basis}
\label{sec:p13-boger-reduction}

Completion establishes that a pattern can serve as a rule. Reduction in
\emph{Prior Analytics} A7 asks which already established rules are redundant.
Boger reads A23's separation of ``completed by means of'' and ``reduced to''
as evidence that the two projects differ. A deduction is used in both projects,
but that common instrument does not identify their aims
\parencite[pp.~207--214]{Boger2004UnderlyingLogic}.

\begin{hbdefinition}{def:p13-boger-reduced-system}{Reduced categorical basis}
Write \(\Gamma\vdash_{\mathrm{red}}c\) for the same finite-chain bookkeeping
with only Barbara and Celarent as admitted \emph{two-premise} rules.
CR1--3, repetition, the stated opposition policy, and indirect deduction
remain available. The restriction removes twelve syllogism rules, not all
auxiliary reasoning and not every existence assumption.
\end{hbdefinition}

\begin{hbtheorem}{thm:p13-boger-reduced-basis}{proved=annex:p13-boger-reduced-basis}{The reduced basis supplies all fourteen syllogism rules}
Under \cref{def:p13-boger-deduction-domain}, every displayed SR1--14 rule
has a finite derivation in the reduced system. Consequently, any finite
full-system deduction can be expanded to a finite reduced-system deduction
with the same undischarged premises and conclusion.
\end{hbtheorem}

The annex gives the twelve nonbasis cases and the finite replacement argument.
In particular, Festino now proceeds indirectly through Celarent; the earlier
direct completion through Ferio is not its reduced-basis proof. Darii and
Ferio also require indirect deductions. The source's table~35 identifies their
routes through Camestres and Cesare, respectively; expanding those routes
exposes Celarent rather than leaving a redundant rule hidden in the proof.

Although Boger describes this project as establishing independence of a small
rule set, the construction just stated establishes \emph{sufficiency} and
redundancy of the removed syllogism rules. A proof that Barbara is not obtainable
without Barbara, or that Celarent is not obtainable without Celarent, would
need a separate nondeducibility argument relative to \emph{all} retained
auxiliaries. No such independence theorem is asserted here merely from the
twelve replacement deductions.

\section{Analysis: transformations and their limits}
\label{sec:p13-boger-analysis}

Analysis in A45 transforms a syllogism in one figure into one in another
through conversion and premise transposition, while retaining the kind of
problem proved: universal attributive, universal privative, partial
attributive, or partial privative. It is neither completion of a validity claim
nor removal of a redundant deduction rule. The figures have no privileged
order for analysis, unlike the first figure's role in completion and reduction
\parencite[pp.~214--216, 241]{Boger2004UnderlyingLogic}.

The following reproduces the \emph{destinations} in Boger's table~37 in new
prose, explicitly marking the source's noted omissions and the transitions
that need premise transposition and conclusion conversion.

\begin{tabularx}{\textwidth}{@{}l X@{}}
\toprule
Starting pattern & Cross-figure destinations in the source's analysis account\\
\midrule
Barbara & No analysis to another figure\\
Celarent & Cesare; Camestres (not treated in the ancient passage)\\
Darii & Datisi; Disamis (not treated in the ancient passage)\\
Ferio & Festino; Ferison\\
Cesare & Celarent\\
Camestres & Celarent, with transposition and conclusion conversion\\
Festino & Ferio; Ferison\\
Baroco & No analysis to another figure\\
Darapti & Darii\\
Felapton & Ferio; Festino\\
Disamis & Darii, with transposition and conclusion conversion\\
Datisi & Darii\\
Bocardo & No analysis to another figure\\
Ferison & Ferio; Festino\\
\bottomrule
\end{tabularx}

The caveat is directedness. Conversion per accidens can replace a universal
premise by a particular converse, losing information while keeping the
intended partial conclusion obtainable. It cannot strengthen that particular
back to the original universal. Therefore an analysis link is not automatically
an equivalence of premise sets, and sharing a conclusion kind alone does not
authorize any desired cross-figure transformation.

\begin{hbtheorem}{thm:p13-boger-analysis-links}{proved=annex:p13-boger-analysis-links}{The seventeen displayed positive analysis links}
Under the ordinary categorical policy, the seventeen positive destinations
in the table can each be realized by conversion and, where indicated,
premise transposition and simple conclusion conversion. Each resulting
syllogism establishes the starting conclusion in its prescribed orientation.
This claim does not assert the negative ``no analysis'' classifications as
an independently proved exhaustive transformation theorem.
\end{hbtheorem}

For example, Darapti has \(PaM,SaM\) and aims at \(PiS\). CR2 changes
\(SaM\) to \(MiS\), giving the Darii pair \(PaM,MiS\). This is a
one-way passage: the existence in \(SaM\) licenses the particular converse,
but \(MiS\) does not license \(SaM\). In the source's account Baroco and
Bocardo retain a nonconvertible \(o\) premise, while Barbara's universal
attributive problem has no counterpart in another figure. These limitations
are why ``all rules are reducible'' does not mean ``every argument can be
analyzed into any figure.''

\section{Finite chains, assumptions, and necessary consequence}
\label{sec:p13-boger-chain-consequence}

Boger extracts formal deducibility from Aristotle's discussion of common
terms, finite reasoning, and rule-governed chains. The narrow two-premise,
three-term syllogism is an elementary link; an extended syllogistic
argumentation is not restricted to two original premises. A conclusion already
reached can serve as a premise in a later link. The middle connects the extreme
terms, rather than leaving separate predications about each disconnected
\parencite[pp.~216--227]{Boger2004UnderlyingLogic}.

For the chain in the source's discussion of \emph{Posterior Analytics} A25,
take the four universal affirmative premises \(AaB,BaC,CaD,DaE\).
Barbara on the first two gives \(AaC\); Barbara with \(CaD\) gives
\(AaD\); Barbara with \(DaE\) gives \(AaE\). The last sentence has four
original premises, although each inferential link uses only two available
sentences and three term places. The source contrasts the longer path through
\(B,C,D\) with a shorter path through \(F,G\) to discuss explanatory priority;
proof length alone does not supply truth or scientific first principles.

Indirect reasoning has a different dependency structure from simply adding a
false statement permanently. To establish \(NoX\) in Baroco, \(NaX\) is
temporarily added, Barbara establishes \(MaX\), and the actual \(MoX\)
makes the subdeduction impossible. The contradictory assumption, not the
actual premise, is discharged. The resulting conclusion depends only on the
original premises. Boger's A23/B14 passages expressly distinguish the
categorical deduction of the impossible result from proof of the original
problem by the hypothesis \parencite[pp.~227--229]{Boger2004UnderlyingLogic}.

The published schematic table~40 repeats \(MaN\) as the result of its
Barbara application although the premises \(MaN,NaX\) yield \(MaX\).
The surrounding sentence, the subsequent contradiction, and the earlier
table~19 fix the intended result. The active reader deductions use \(MaX\);
they do not present that source-table slip as a valid new inference.

For necessary consequence Boger assembles passages about a result's not being
otherwise, agreement of truths, and inability to deduce a falsehood from truths.
His substitution formulation says: every admissible term substitution making
the premises true makes the conclusion true. It is a historical reconstruction
of the notion, not a claim that Aristotle wrote a modern model-theoretic
definition or that substitution and reinterpretation are identical concepts
\parencite[pp.~230--235, 242--243]{Boger2004UnderlyingLogic}.

The source's table~43 offers an indirect Barbara argument: from \(AaB,BaC\),
assume \(AoC\); Baroco applied to \(AaB,AoC\) yields \(BoC\), conflicting
with \(BaC\); discharge \(AoC\) to conclude \(AaC\). This is legitimate
\emph{inside} the already established full rule inventory. Because Baroco's
completion itself used Barbara, the table is not a noncircular foundational
proof of Barbara independent of its earlier admission.

Boger also reads \emph{Metaphysics} IV as recognizing consequence from
contradiction, while his note~104 records Priest's disagreement with that
reading. The categorical replacement theorem above neither requires adoption
of every sentence as a categorical consequence of inconsistent premises nor
settles that interpretive dispute. An agreed incompatible pair is used for
assumption discharge in the expressly stated local system.

\begin{proofboundary}{def:p13-boger-deduction-adequacy-boundary}{What these constructions establish}
These deductions establish the ten completions, the smaller sufficient
two-premise rule basis with a finite expansion procedure, the seventeen positive
analysis links, and the stated contrasted-instance and weakening implications.
They do not prove a general completeness theorem for arbitrary categorical
premise sets, a decision algorithm with a termination proof, independence of
the retained basis, or equivalence with an unrestricted first-order logic.
Boger reads A23 and A30 as expressing a completeness aspiration, but the
finite categorical arguments here are not silently promoted to that global
metatheorem. Nor is his historical interpretation an independently admitted
Aristotle text edition.
\end{proofboundary}

\begin{proofboundary}{def:p13-boger-deduction-source-boundary}{Exact source-treatment limits}
The source's surrounding prose and the forty-eight-pair classification are
available, but the held scan cuts the right columns of catalogue tables~26--28
(printed pp.~198--200) and the right-hand entries of reduction table~36
(pp.~212--213). These columns are also cut in the original embedded scan images,
not merely lost by English OCR. The deductions supplied here are complete
authored constructions based on the legible rule statements and argument,
not fabricated transcriptions of the clipped cells. Exact recovery/collation
of those cells and the Greek quotations remains a source-fidelity duty.
Johnson's later modal apparatus and the full two-chapter Greek allocation
remain distinct from this categorical continuation.
\end{proofboundary}
